\documentclass[IJ]{cesj}
\usepackage{graphicx}
\usepackage{listings}

\author{Ga\"etan Lehmann}
\institute{Biologie du d\'eveloppement et de la reproduction, INRA de Jouy-en-Josas}

\title{ModulusImageFilter}
\abstract{ModulusImageFilter is a simple image filter to compute the modulus (x \% dividend) pixel-wise.}
\keyword{Modulus UnaryFunctorImageFilter}
\year{2005}

\leftmark{Ga\"etan Lehmann}
\rightmark{Ga\"etan Lehmann}

\begin{document}
\lstset{language=c++}
\maketitle
% \tableofcontents

\section{Description}
ModulusImageFilter is a simple image filter to compute the modulus (x \% dividend) pixel-wise.

\section{Implementation}
ModulusImageFilter is based on UnaryFunctorImageFilter and use the c++ modulus operator (\%).

\section{Usage}
As usual, user has to include the header:
\begin{lstlisting}
#include "itkModulusImageFilter.h"
\end{lstlisting}
and declare the type, instanciate the filter, and set the input image
\begin{lstlisting}
typedef itk::ModulusImageFilter< InputType, OuputType >
  ModulusType;
ModulusType::Pointer modulus = ModulusType::New();
modulus->SetInput( filter->GetOutput() );
\end{lstlisting}
The dividend can be set with the SetDividend() method, and read with GetDividend(). The default value is 5.
\begin{lstlisting}
modulus->SetDividend( 8 );
\end{lstlisting}


\section{Example}
ModulusImageFilter is used here to highlight the level curves of the image produced by the DanielssonDistanceMapImageFilter filter.

% \begin{figure}[!hbp]
% \centering
% \includegraphics[width=0.5\textwidth]{Spots.eps}
% \caption{The input image.}
% \end{figure}
% 
% \begin{figure}
% \centering
% \includegraphics[width=0.5\textwidth]{distance.eps}
% \caption{The (rescaled) distance map image.}
% \end{figure}
% 
% \begin{figure}
% \centering
% \includegraphics[width=0.5\textwidth]{modulus.eps}
% \caption{The (rescaled) output image.}
% \end{figure}

\end{document}
